\chapter{Fazit und Ausblick}
Im Rahmen dieser Masterarbeit wurde die Grundlage für die modellbasierte Funktionsentwicklung im Bereich der Akkuprodukte bei STIHL geschaffen. Durch eine sorgfältige Vorbereitung des Matlab-Modells, insbesondere hinsichtlich der Modellstruktur und der automatisierten Code-Generierung, konnte ein solides Fundament für die Arbeit gelegt werden. Mithilfe von Variant-Subsystemen wurde die bestehende modulare Firmware-Architektur erhalten und erweitert.\\
Eine umfassende Analyse des Lastenhefts sowie des bestehenden C-Codes bildete die Basis für die Modellierung und Implementierung der Maschinenschutzfunktionen. Mithilfe der Simulink-Test-Toolbox wurden die einzelnen Funktionen in einer isolierten Testumgebung unter Verwendung selbst definierter Testabläufe automatisiert geprüft. Durch die Verknüpfung der Anforderungen aus dem Lastenheft mit den definierten Testabläufen konnten die implementierten Funktionen automatisch verifiziert werden. \\
Der C-Code für die Maschinenschutzfunktionen wurden mittels automatisierter Codegenerierung durch den Embedded-Coder erstellt und in die bestehende Firmware integriert. Über eine Konfigurationsoption in der Header-Datei der jeweiligen Maschine ist es möglich, während der Erstellung des Maschinencodes zwischen den automatisch generierten Funktionen und den vorhandenen Funktionen der Basissoftware zu wechseln. Bei aktivierter Autocode-Option werden die entsprechenden Basissoftware-Funktionen deaktiviert und während des Kompilierens des gesamten Codes nicht berücksichtigt.\\
Im Rahmen nachfolgender Systemtests wurden die modellbasiert entwickelten Maschinenschutzfunktionen erfolgreich an einer akkubetriebenen Motorsäge geprüft und verifiziert. Eine abschließende Analyse des Speicherbedarfs und der Laufzeit bei aktiviertem und deaktiviertem Autocode veranschaulichte die Unterschiede zwischen dem Autocode und den manuell codierten Funktionen. Bei aktivierten Autocode-Funktionen erhöht sich der Speicherbedarf, primär bedingt durch die Schnittstelle zwischen der Basissoftware und dem Autocode, die für die Bereitstellung der Ein- und Ausgangssignale verantwortlich ist. Die Speichererhöhung bezogen auf den ROM beläuft sich lediglich auf \unit[$0,37$]{\%}, wenn der zusätzliche Speicherbedarf der Schnittstelle berücksichtigt wird, und auf \unit[$0,107$]{\%}, wenn ausschließlich die Funktionalität in Relation zum gesamten verfügbaren Speicher des Steuergeräts betrachtet wird. Auch die Laufzeit erhöht sich geringfügig aufgrund der zusätzlichen Abarbeitungszeit der Schnittstelle, was jedoch keinen signifikanten Einfluss auf die Gesamtperformance hat. Mit der zunehmenden Auslagerung von Funktionen aus der Basissoftware in den Playground relativiert sich der zusätzliche Speicherbedarf sowie die verlängerte Laufzeit der Schnittstelle.\\
Der modellbasierte Entwicklungsansatz führt im Vergleich zur manuellen Programmierung zu einer geringfügigen Erhöhung des Speicherbedarfs und der Laufzeit. Jedoch überwiegen die Vorteile dieses Ansatzes. Die strukturierte Implementierung der Funktionen auf Modellebene verbessert deren Nachvollziehbarkeit und Verständlichkeit erheblich. Dadurch können präzise Anpassungen einzelner Funktionen effizient und mit geringem Aufwand realisiert werden.

Zur abschließenden Prüfung und Validierung der in dieser Arbeit entwickelten Maschinenschutzfunktionen auf der realen Hardware sollte die Maschine mit aktiviertem Playground den standardisierten Prüfablauf der Testing-Gruppe durchlaufen.\\ 
Die Fortführung des Themas der modellbasierten Funktionsentwicklung ermöglicht es, weitere Funktionen aus der bestehenden Basissoftware in das Matlab-Modell zu übertragen. Hierbei kann die erarbeitete Modellstruktur sowie die konfigurierte Autocodegenerierung verwendet werden. Darüber hinaus besteht die Möglichkeit, die bestehende Schnittstelle zwischen der Basissoftware und dem Autocode weiter zu verwenden, Anpassungen vorzunehmen oder bei Bedarf zu erweitern.\\
Zusätzliche Untersuchungen könnten sich auf die Speicherbelegung in Bezug auf die Autocode-Einstellungen und die Modellierung konzentrieren. Derzeit erfolgt die Signalübergabe zwischen den Maschinenschutzfunktionen im Autocode durch den Zugriff auf globale Variablen. 
Zur Optimierung des Speicherbedarfs sollte dieses Verfahren überprüft und gegebenenfalls angepasst werden, um den Speicherverbrauch weiter zu minimieren. Dadurch bestehen weiterhin Möglichkeiten, die Leistungsfähigkeit und Effizienz des Autocodes zu verbessern.\\
Parallel zu dieser Arbeit wird ein physikalisches Modell in Matlab entwickelt, welches das Gesamtverhalten von Akkugeräten simuliert. Dieses Modell soll sowohl die physikalischen Eigenschaften als auch die Softwarefunktionen enthalten. Die im Rahmen dieser Arbeit implementierten Maschinenschutzfunktionen, die in einem Referenced-Subsystem abgespeichert sind, lassen sich mühelos in das physikalische Modell integrieren. Dabei ist lediglich die Übergabe der Eingangssignale sowie die Verarbeitung der Ausgangssignale zu berücksichtigen.\\
Zusammengefasst wird deutlich, dass der modellbasierte Ansatz auch künftig vielfältige Einsatzmöglichkeiten bei STIHL im Bereich der Akkugeräte bietet und zusätzliche Vorteile bei der Produktentwicklung generiert.

% Im Rahmen dieses Projekts wurden erste Entwicklungsschritte für die Entwicklung eines "`schlangenförmigen"' Multigelenk-Roboters, der im Bereich der Medizintechnik unter begrenzten Platzverhältnissen eingesetzt werden kann, ausgeführt. Bei bestimmten Untersuchungen und chirurgischen Eingriffen ist es von Vorteil, einen "`schlangenartig"' beweglichen Roboterarm einzusetzen, der seine Form an die Gegebenheiten innerhalb des menschlichen Körpers anpassen kann. Durch die gezielte Verformung des Roboterarms in allen drei Dimensionen, werden Verletzungen, die während eines Eingriffs möglicherweise verursacht werden minimiert und es können Stellen innerhalb des menschlichen Körpers erreicht werden, die mit starren Strukturen nicht zugänglich sind.\\
% Basierend auf den aus dem ausgewählten Anwendungsgebiet, der Darmspiegelung, abgeleiteten Anforderungen wurden Lösungsansätze für einen Multigelenk-Roboter abgeleitet. Diese Lösungsansätze wurden mit verschiedenen Bewertungsmethoden zu Konzeptvarianten kombiniert. Die ausgewählte Konzeptvariante orientiert sich an einer kardanischen Aufhängung. Dies ermöglicht eine Bewegung mit zwei rotatorischen Freiheitsgraden an einem Gelenk. Durch das aneinanderreihen  mehrerer solcher Gelenke kann eine "`schlangenförmige"' Bewegung generiert werden.\\
% Im Rahmen einer Machbarkeitsstudie wurde untersucht, ob die geforderte Bewegung durch den Einsatz von SMA-Drähten als Aktuatoren ausgeführt werden kann. Dafür wurde das gewählte Ansteuerkonzept in einem Funktionsmodell, welches in mehreren Anpassungsschritten optimiert wurde, getestet. Das antagonistische Anordnen der SMA-Drähte ermöglicht dabei eine rotatorische Bewegung in zwei Richtungen. Im letzten Schritt wurden die aus dem Funktionsmodell abgeleiteten Erkenntnisse und Winkelauslenkungen, die ein Gelenk ausführen kann auf die geforderte Zielbaugröße umgerechnet und in die gewählte Konzeptvariante integriert. Durch die um 120° versetzten Anordnung der Ansteuermechanismen ist eine rotatorische Bewegung in zwei Dimensionen an jedem Gelenk möglich. Die Berechnungen ergeben, dass die geforderte Bewegung an einem Gelenk ausgeführt werden kann. Somit wurde mit den durchgeführten Untersuchungen gezeigt, dass ein Multigelenk-Roboter auf Basis des ausgewählten Konzepts und des gewählten Ansteuermechanismus entwickelt werden kann.

% In einem nächsten Entwicklungsschritt müsste die ausgeführte Berechnungen für die Kombination aus dem ausgewählten Gelenkmechanismus (kardanische Aufhängung) und den drei um 120° versetzten Ansteuermechanismen (3 Hebelkaskaden) in einem neuen Funktionsmodell getestet werden. Auch hier ist es zunächst möglich, dass neue Funktionsmodell in einem vergrößerten Maßstab aufzubauen. In einem weiteren möglichen Entwicklungsschritt kann der Ansteuermechanismus in der Originalgröße ausgeführt werden, um zu überprüfen, ob die im Funktionsmodell ermittelten Erkenntnisse und Rotationswinkel auch in der Zielgröße übereinstimmen. Bevor hier ein reales Modell aufgebaut wird, können die entstehenden Rotationswinkel für die minimierten Abmessungen zunächst in der Simulation mit Matlab-Simulink getestet werden. Ein Testen, bei dem die SMA-Drähte gegen ein Gegengewicht ziehen, ist mit dem bereits erstellten Simulationsmodell durch Anpassen der Eingangsgrößen möglich. Für das Simulieren einer antagonistischen Drahtanordnung müsste das Simulationsmodell auf eine antagonistische Drahtanordnung abgeändert werden.
